\documentclass[../../main.tex]{subfiles}
\begin{document}


{\bf [解]}
題意より
\begin{align}
0 \le \alpha < \pi, \quad 0 < \beta < \frac{1}{2} \label{eq:1} \\
0 \le x \le 1, \quad y \le \cos(\pi\beta) \label{eq:2}
\end{align}
である．
\cref{eq:1}より $0 < \alpha\beta < \dfrac{\pi}{2}, \, 0 < \pi\beta < \dfrac{\pi}{2}$ だから, 
\begin{align}
 \cos(\pi\beta) > 0, \, \sin(\pi\beta) > \sin(\alpha\beta) > 0 \label{eq:3}
\end{align}
となる．
ここで
\begin{align}
 A(x,y,\alpha,\beta) = y \sin(\alpha\beta) + x \sin\{(\pi-\alpha)\beta\} \label{eq:4}
\end{align}
とおく．$y$以外を固定して$y$を動かすと\cref{eq:2,eq:3}より$A$は$y = \cos(\pi\beta)$ で最大値をとる．これを$B(x,\alpha,\beta)$とおく：
\begin{align}
  B(x,\alpha,\beta) = \cos(\pi\beta) \sin(\alpha\beta) + x \sin\{(\pi-\alpha)\beta\}. 
\end{align}
次に$x, \beta$を固定して$\alpha$を動かす．$B$の$\alpha$に関する一階微分は
\begin{align} 
\frac{dB}{d\alpha} 
   &= -\beta [ \cos(\pi\beta) \cos(\alpha\beta) + x \cos\{(\pi-\alpha)\beta\} ] \\
   &\le 0 \quad (\because \text{\cref{eq:1,eq:2}}) 
\end{align}
であるから， $B$ は $\alpha$ の単調減少関数で\cref{eq:1}より $\alpha = 0$ で最大値をとる．これを $C(x,\beta)$ とおく：
\begin{align} 
 C(x,\beta) 
  &= x \sin\pi\beta 
\end{align}
ここで\cref{eq:2,eq:3}よりこれは$x=1$で最大値$\sin\pi\beta$をとる．
以上から
\begin{align}
 A \le \sin\pi\beta \label{eq:5}
\end{align}
である．

従って，与式$\sin\pi\beta \le A$が成立するのは\cref{eq:5}で等号が成立する時のみで，この時$x=1$である．

\newpage

\end{document}
